% !TeX root = ../深度学习笔记.tex
% 由主文档通过 \input 载入；请编译主文档。
\setcounter{chapter}{10} % 保留原书章号。
\chapter{残差网络与视觉任务}
\label{chap:dl-vision}
加深网络可以组合更多局部变换，但也增加优化难度。残差连接为信号提供直接路径；分类之外的视觉任务则要求网络保留或恢复空间信息，因而需要不同输出结构。

\section{残差连接与深层网络}
\emph{残差连接（residual connection）}把输出写成 $\mathbf y=\mathbf x+F_\theta(\mathbf x)$，两项必须同形。在可微点 $D\mathbf y=\mathbf I+DF_\theta$，反向梯度包含直接传递的 $\nabla_{\mathbf y}\ell$ 与经过残差分支的贡献。即使分支接近零，块仍可近似恒等映射，使“保留输入”成为容易表达的状态。

恒等项不保证任意深网络的梯度都不消失：分支导数仍可能相消，多个块的雅可比仍需连乘。若通道数或空间尺度变化，可写成 $\mathbf y=P\mathbf x+F_\theta(\mathbf x)$，其中 $P$ 是相容投影，常用带适当步幅的 $1\times1$ 卷积实现。此时直接路径导数是 $P$，不再是同维恒等矩阵。

\section{模块化设计与多尺度表示}
$1\times1$ 卷积在每个位置混合通道，不直接扩展空间感受野。瓶颈块先减少通道，再进行空间卷积，最后恢复通道，可降低中间计算。下采样层扩大后续单元对应的输入范围，却降低定位精度；多尺度结构把不同分辨率特征结合，以兼顾局部细节和较大范围信息。

\emph{深度可分离卷积（depthwise separable convolution）}先对每个输入通道独立作空间卷积，再用 $1\times1$ 卷积混合通道。忽略偏置、设通道倍率为一，参数量从 $k^2C_{\mathrm{in}}C_{\mathrm{out}}$ 降为 $k^2C_{\mathrm{in}}+C_{\mathrm{in}}C_{\mathrm{out}}$，但结构约束也不同。理论乘加量下降不保证同比加速，小算子调度和内存访问可能占主导。

\section{目标检测的输出与评价}
\emph{目标检测（object detection）}输出一组对象，每个对象包含类别和边界框，可用 $(x_{\min},y_{\min},x_{\max},y_{\max})$ 表示。对面积有限且并集非空的两个区域 $A,B$，\emph{交并比（intersection over union, IoU）}为 $|A\cap B|/|A\cup B|$。分类损失评价类别，位置回归损失评价边界；训练前必须说明哪个预测与哪个真实对象匹配，未匹配预测如何处理。

候选框方法生成多个候选并评分，集合式方法直接学习有限预测集合及匹配。\emph{非极大值抑制（non-maximum suppression, NMS）}依次保留高分候选，删除与其重叠过大的同类候选，是减少重复预测的后处理。平均精度依据按分数排序并匹配后的精确率--召回率曲线计算；不同 IoU 阈值、插值和类别汇总规则会给出不同数值，比较时应固定完整协议。

\section{语义分割与编码器--解码器}
\emph{语义分割（semantic segmentation）}为每个像素赋予类别；\emph{实例分割（instance segmentation）}还要区分同类的不同对象。编码器逐步降低分辨率以提取语义，解码器恢复空间输出，跳跃连接把较高分辨率特征送到对应解码层。像素交叉熵可在有效标注位置平均，类别权重会改变不同像素的训练贡献。

上采样可用固定插值或可训练的\emph{转置卷积（transposed convolution）}。将普通卷积写成线性映射 $\mathbf y=\mathbf C\mathbf x$，转置卷积对应 $\mathbf C^\mathsf T$，一般不是 $\mathbf C^{-1}$。在一维、膨胀率一的常见设置下，其输出长为 $(H_{\mathrm{in}}-1)s-2p_0+k+o$，其中输出补偿 $0\leq o<s$ 选择相容形状。输出与标签须逐像素对齐，否则损失可能在错误位置比较。

\section{例题与习题}
\begin{example}[区域重叠]
两个边界框分别为 $A=[0,2]\times[0,2]$ 与 $B=[1,3]\times[0,2]$。求 IoU，并解释其含义。
\end{example}
\begin{solve}
两框面积各为 $4$，交集面积为 $2$，并集面积为 $4+4-2=6$，故 IoU 为 $1/3$。它衡量位置重叠，不说明类别是否正确；类别预测正确而位置偏离仍可能在检测评价中算作失败。
\end{solve}
\begin{enumerate}
\item 设计把 $32\times32\times16$ 特征变为 $16\times16\times32$ 的残差块，写明两条分支形状。
\item 对 $C_{\mathrm{in}}=C_{\mathrm{out}}=32,k=3$ 比较普通卷积和深度可分离卷积的无偏置参数量。
\item 用合成矩形图像建立分割实验，比较插值和转置卷积上采样，并分别报告像素指标、区域指标与推断成本。
\end{enumerate}

